%%%%%%%%%%%%%%%%%%%%%%%%%%%%%%%%%%%%%%%%%%%%%%%%%%%%%%%%%%%%%%%%%%%%%%%%
% Plantilla TFG/TFM
% Universidad de Murcia. Facultad de Informática
% Realizado por: José Manuel Requena Plens
% Modificado: Pablo José Rocamora Zamora
% Contacto: pablojoserocamora@gmail.com
%%%%%%%%%%%%%%%%%%%%%%%%%%%%%%%%%%%%%%%%%%%%%%%%%%%%%%%%%%%%%%%%%%%%%%%%


\chapter{Estado del arte}

Este capítulo recorre las bases teóricas sobre las que se construye el sistema propuesto. Se parte de la arquitectura Transformer y el mecanismo de atención \cite{vaswani2017attention} que hizo posibles los LLMs actuales, para llegar a una de sus limitaciones más críticas: la alucinación. Frente a este problema, la arquitectura RAG inyecta documentos externos relevantes en cada consulta para que el modelo los use como base factual al generar su respuesta, lo que desplaza parte de la responsabilidad factual al componente de recuperación. Por ello, el capítulo también traza la evolución de los sistemas de recuperación de información, desde los modelos léxicos clásicos hasta la búsqueda semántica por vectores, cuya comparativa constituye uno de los ejes experimentales de este trabajo. Finalmente, se abordan los modelos generadores empleados y los métodos de evaluación utilizados para medir la calidad de sus respuestas.

\section{Paso del modelo estadístico al mecanismo de atención}

Durante décadas, los sistemas de procesamiento del lenguaje natural se apoyaron en modelos estadísticos que trataban el texto como una secuencia de símbolos con probabilidades de co-ocurrencia. Arquitecturas como las redes neuronales recurrentes (RNN) y sus variantes con memoria a largo plazo (LSTM) procesaban la información de forma estrictamente secuencial, imponiendo dos límites estructurales severos: la incapacidad de paralelizar el cómputo durante el entrenamiento y la dificultad de capturar dependencias entre palabras distantes en el texto, ya que la información sobre el inicio de una oración tendía a diluirse a medida que el modelo avanzaba hacia su final.

El punto de inflexión llega en el 2017 con la proposición de una arquitectura radicalmente diferente, el \textit{Transformer}, un modelo que prescinde por completo de la recurrencia y las convoluciones, apoyándose en su lugar de forma exclusiva en un mecanismo denominado \textit{self-attention} \cite{vaswani2017attention}.

El mecanismo de atención de la arquitectura Transformer permite que cada elemento de una secuencia atienda directamente a cualquier otro, evitando la pérdida de información de los modelos tradicionales al comprimir el contexto previo en un único vector de estado. A nivel formal, el sistema calcula puntuaciones de relevancia entre un conjunto de consultas (\textit{queries}), claves (\textit{keys}) y valores (\textit{values}), lo que genera una representación contextual ponderada de forma paralela. Esta capacidad se potencia mediante el enfoque \textit{multi-head attention}, el cual dota al modelo de la habilidad para proyectar y capturar simultáneamente múltiples tipos de relaciones lingüísticas entre los términos de la secuencia \cite{vaswani2017attention}.

Al eliminar la secuencialidad del cómputo, el Transformer podía aprovechar de mucho mejor el paralelismo de las GPUs modernas, lo que abrió la puerta al entrenamiento de modelos con decenas y centenares de miles de millones de parámetros. Sobre esta arquitectura se construyeron los modelos que hoy conocemos como LLMs, entre ellos los modelos de las familias Llama y Phi que se evalúan en este trabajo.

Sin embargo, la capacidad de generar texto fluido y contextualmente coherente no implica capacidad de razonar sobre la veracidad de los hechos. Esta distinción, que la arquitectura Transformer no resuelve por diseño, constituye el problema central abordado en la sección siguiente.

\section{Alucinación en modelos de lenguaje}

Los LLMs tienen una tendencia documentada a producir texto factualmente incorrecto con apariencia de certeza. En el ámbito del procesamiento del lenguaje natural, se denomina alucinación a la generación de contenido fluido y sintácticamente correcto que, sin embargo, es factualmente inexacto o carece de respaldo en fuentes externas \cite{alansari2025hallucination}. Lo que hace especialmente problemático este fenómeno es que el texto alucinado no se distingue superficialmente del texto verídico, lo que dificulta tanto su detección automática como la manual.

La razón estructural de este comportamiento reside en el objetivo de entrenamiento. La mayoría de los LLMs se preentrenan mediante la técnica de estimación de máxima verosimilitud, basada en la predicción del siguiente \textit{token}. Bajo este paradigma, el modelo aprende a asignar la mayor probabilidad posible a la continuación observada en los datos de entrenamiento \cite{welleck2019neural}, maximizando la verosimilitud estadística del corpus en lugar de su correspondencia factual con la realidad. Como resultado, un modelo que predice con éxito el siguiente token captura patrones lingüísticos y co-ocurrencias semánticas precisas, pero carece de mecanismos intrínsecos que le obliguen a distinguir entre lo que es estrictamente verdadero y lo que es simplemente plausible \cite{alansari2025hallucination}.

Existe además una razón estadística por la que los LLMs alucinan con independencia de su arquitectura o de la calidad de sus datos. \citet{kalai2023calibrated} demostraron que, incluso partiendo de datos de entrenamiento perfectamente verídicos, existe una tasa mínima de alucinación que ningún modelo puede evitar. El motivo es que hay hechos cuya veracidad simplemente no puede aprenderse del corpus: si un hecho concreto aparece solo una vez en los datos de entrenamiento, el modelo no tiene información suficiente para saber si es verdadero o inventado. En dominios como el \textit{Fact-Checking}, donde abundan fechas, cifras y nombres propios poco frecuentes, este problema es especialmente acusado.

\citet{bender2021parrots} denominaron a los LLMs \textit{loros estocásticos} (\textit{stochastic parrots}): sistemas que ensamblan secuencias de formas lingüísticas siguiendo información probabilística sobre cómo se combinan, pero sin ninguna referencia al significado. La fluidez del texto generado es una consecuencia del aprendizaje estadístico, no un indicador de comprensión ni de veracidad. Un LLM operando de forma aislada, sin acceso a fuentes externas verificables, no es por tanto un componente fiable para verificación de información, lo que motiva la adopción de la arquitectura RAG descrita en el capítulo de análisis.

\section{Recuperación de información: de la búsqueda por palabras clave a los vectores semánticos}

La recuperación de información (\textit{Information Retrieval}, IR) estudia cómo localizar, dentro de una colección de documentos, aquellos que son relevantes para una consulta. Su evolución reciente puede dividirse en dos etapas bien diferenciadas: el paradigma léxico, cuyo exponente principal es BM25, y el paradigma semántico basado en representaciones vectoriales densas.

\subsection{BM25 y la búsqueda léxica}

El modelo BM25 (\textit{Best Matching 25}), desarrollado entre otros por \textit{Stephen E. Robertson} en el contexto del sistema Okapi de la Universidad de Londres durante los años ochenta y noventa, es la función de puntuación probabilística de referencia en recuperación de información \cite{robertson2009bm25}. Estima la relevancia de un documento para una consulta combinando dos señales: la frecuencia de los términos de la consulta en el documento (\textit{term frequency}, TF) y la rareza de esos términos en el conjunto total de documentos (\textit{inverse document frequency}, IDF), con un mecanismo de normalización por longitud del documento que corrige el sesgo hacia textos más extensos \cite{robertson2009bm25}. Su eficiencia computacional, basada en índices invertidos, lo mantuvo durante décadas como  referencia en motores de búsqueda industriales. Elasticsearch lo adopta como algoritmo  predeterminado, razón por la que resulta el punto de partida natural de la comparativa  experimental de este trabajo.

BM25 funciona comparando coincidencias léxicas exactas, por lo que si una consulta usa una palabra como ``fallecimiento'' y el documento que contiene la respuesta emplea ``muerte'' o ``deceso'', BM25 no detecta ninguna relación y el documento no se recupera \cite{karpukhin2020dpr}. En \textit{Fact-Checking}, donde una misma noticia puede formularse de formas muy distintas a como se plantea la pregunta de verificación, esta rigidez léxica es una barrera real que se podrá ver más adelante en los análisis y resultados del trabajo.

\subsection{Recuperación semántica mediante \textit{embeddings}}

\citet{karpukhin2020dpr} demostraron que aplicar la arquitectura Transformer a la recuperación permite proyectar consultas y documentos a un espacio vectorial denso donde la similitud geométrica refleja la relevancia semántica. Dos textos sobre el mismo concepto quedan próximos en ese espacio aunque no compartan ninguna palabra. Los experimentos de los autores mostraron que este enfoque superaba a BM25 en entre 9 y 19 puntos porcentuales en precisión de recuperación \textit{top-20} sobre conjuntos de datos de preguntas abiertas \cite{karpukhin2020dpr}. La representación resultante, el \textit{embedding}, es un vector de cientos o miles de dimensiones que codifica el significado del texto; la recuperación se reduce entonces a una búsqueda del vecino más próximo, operación para la que existen implementaciones eficientes incluso a escala de cientos de miles de documentos.

A raíz de esto han surgido distintas familias de modelos de \textit{embeddings}. \citet{reimers2019sbert} propusieron Sentence-BERT (SBERT), que adapta BERT mediante redes siamesas y de trillizos para obtener representaciones semánticamente significativas comparables por similitud coseno, reduciendo el tiempo de búsqueda sobre colecciones de diez mil oraciones de 65 horas a apenas cinco segundos respecto al uso directo de BERT. En cuanto a modelos multilingües, Microsoft presentó los modelos mE5, entrenados sobre más de mil millones de pares de texto en diversas lenguas, alcanzando un rendimiento comparable al de los mejores modelos monolingües en inglés \cite{wang2024multilingual}.

Para este TFG se ha elegido utilizar BGE-M3, presentado por \citet{chen2024bge}, que integra en un único modelo las dos modalidades de recuperación descritas en este capítulo. Soporta más de cien idiomas, lo que lo hace adecuado para corpus en español; unifica recuperación densa, dispersa y multi-vector; y procesa entradas de hasta 8.192 tokens, cubriendo sin dificultad la extensión habitual de una noticia. Disponer de ambas estrategias en el mismo modelo permite comparar la búsqueda léxica y la semántica manteniendo constantes el resto de componentes del sistema, lo que proporciona una medida directa del impacto de cada estrategia de recuperación sobre la veracidad final de las respuestas.

\section{Generación Aumentada por Recuperación}

La arquitectura RAG (\textit{Retrieval-Augmented Generation}) fue formalizada por \citet{lewis2020rag} en 2020 como un marco que combina la memoria paramétrica de un modelo generativo preentrenado con una memoria no paramétrica: un índice denso de vectores consultable mediante un recuperador neuronal. Los autores demostraron que sistemas RAG de apenas unos cientos de millones de parámetros superaban en tareas de preguntas abiertas a modelos paramétricos cerrados de once mil millones, generando texto más específico y factualmente fundamentado \cite{lewis2020rag}.

En cuanto a la arquitectura, un sistema RAG funciona con tres fases bien diferenciadas \cite{gao2024rag}. La primera es la indexación: los documentos de la base de conocimiento se dividen en fragmentos (\textit{chunks}), se codifican mediante un modelo de \textit{embeddings} y se almacenan en un índice vectorial que permite búsquedas eficientes por similitud. La segunda fase es la recuperación: ante una consulta entrante, el sistema la proyecta al mismo espacio vectorial y extrae los fragmentos más relevantes según una métrica de similitud (por ejemplo el producto escalar o la similitud coseno). La tercera fase es la generación: los fragmentos recuperados se concatenan con la consulta original en el prompt del modelo generativo, que produce una respuesta teniendo un contexto relevante en la consulta en lugar de depender únicamente de su conocimiento paramétrico.

El problema es que los documentos irrelevantes contaminan ese contexto, y la posición de los documentos dentro de él tampoco es indiferente para el modelo. \citet{liu2024lost} demostraron que el rendimiento de los LLMs en tareas de acceso a información dentro de contextos largos se degrada significativamente cuando la información relevante queda en posiciones intermedias, siendo máximo cuando aparece al inicio o al final de la ventana \cite{liu2024lost}. Recuperar más documentos de los necesarios, por tanto, no mejora el sistema: puede empeorarlo, lo cual hace imprescindible el ajuste del \textit{top-k} de recuperación para cada modelo y tarea concreta.

\section{Modelos de lenguaje pequeños como generadores}

\citet{abdin2024phi3} presentaron \textit{Phi-3-mini-4k-instruct}, un modelo de 3.800 millones de parámetros entrenado sobre 3,3 billones de tokens, donde la clave no está en la arquitectura sino en la curación del conjunto de entrenamiento: datos web fuertemente filtrados combinados con datos sintéticos de alta calidad. El resultado es que Phi-3-mini alcanza un 69\% en el benchmark MMLU y 8,38 puntos en MT-bench \cite{abdin2024phi3}, cifras comparables a las de Mixtral 8x7B o GPT-3.5 \cite{abdin2024phi3} con un orden de magnitud más de parámetros, y con un tamaño suficientemente reducido como para ejecutarse en un teléfono móvil \cite{abdin2024phi3}.

Esto va a tener implicaciones directas en este TFG. El \textit{Fact-Checking} con RAG es una tarea estructurada: dado un conjunto de documentos recuperados, el modelo debe determinar si la afirmación es coherente con la evidencia y formular una respuesta ceñida a ese contexto, por lo que requiere seguir instrucciones precisas y razonar sobre el contexto recibido. Los propios autores del informe técnico de Phi-3 señalan que los SLM no están diseñados para recuperación de conocimiento en profundidad desde sus pesos \cite{abdin2024phi3}, en nuestro caso el modelo no va a necesitar almacenar el conocimiento porque lo recibe del recuperador. Desde el punto de vista operativo, ejecutar Phi-3 en local elimina dependencias de API externas, reduce los costes de infraestructura y permite evaluar el verdadero impacto del recuperador sobre la veracidad de las respuestas, sin que el modelo pueda compensar con su conocimiento paramétrico las deficiencias de la recuperación.

No obstante, esta reducción en el tamaño del modelo conlleva una pérdida en la profundidad de razonamiento. Aunque los SLM han demostrado una competencia sorprendente en tareas específicas, los LLMs son superiores en lo que respecta a seguir instrucciones altamente complejas y en la resolución de ambigüedades lógicas. Su inclusión en el marco experimental de este trabajo permite evaluar si la eficiencia operativa de los modelos pequeños implica una degradación crítica en la calidad de las respuestas, o si, por el contrario, la inserción de contexto proporcionada por el RAG logra igualar el rendimiento de ambas escalas. Por ello, los SLMs se usarán en primer lugar para evaluar el rendimiento de los recuperadores, para luego comparar su rendimiento con el de los LLMs utilizando el mismo sistema de recuperación.

\section{Evaluación de respuestas (LLM-as-a-judge)}

Evaluar si una respuesta generada es correcta es un problema difícil. Durante décadas la solución dominante fueron métricas de solapamiento de n-gramas: BLEU para traducción automática y ROUGE para resumen. Su limitación es estructural: comparan palabras, no significados \cite{zheng2023judging}. Una respuesta correcta pero formulada con palabras distintas a las de la referencia recibe una puntuación baja mientras que una respuesta incorrecta pero con las mismas palabras en orden diferente puede obtener una puntuación alta. En tareas de verificación factual, donde lo que se evalúa es precisamente el contenido semántico, estas métricas no son adecuadas.

La alternativa a esto es el paradigma \textit{LLM-as-a-judge}, que consiste en utilizar un LLM para evaluar las respuestas de otros modelos. \citet{zheng2023judging} propusieron y estudiaron sistemáticamente este enfoque, introduciendo los benchmarks MT-bench y Chatbot Arena, y verificaron que un juez como GPT-4 alcanza más de un 80\% de acuerdo con las preferencias humanas, el mismo nivel de concordancia observable entre evaluadores humanos. Esta escalabilidad lo convierte en la mejor alternativa viable cuando se evalúan cientos de respuestas.

Sin embargo, el mismo estudio que lo popularizó identificó sus limitaciones estructurales: sesgo de posición, por el que el modelo favorece la respuesta presentada en primer lugar; sesgo de verbosidad, por el que respuestas más largas puntúan mejor con independencia de su calidad; sesgo de autopreferencia, por el que el juez puntúa mejor los textos generados por modelos próximos al suyo; y capacidad de razonamiento limitada en verificación factual compleja \cite{zheng2023judging}. Investigaciones posteriores han profundizado en el sesgo de autopreferencia: su origen no está en que el modelo reconozca sus propias respuestas, sino en que tiende a favorecer textos con menor perplejidad bajo su distribución, es decir, textos que le resultan más familiares estadísticamente \cite{wataoka2024selfpref}. Esto lo convierte en un sesgo inherente al mecanismo de evaluación, difícil de mitigar solo con ingeniería de prompts.

En este trabajo, el paradigma \textit{LLM-as-a-judge} se adopta como mecanismo de evaluación principal por su capacidad de escalar a los 220 casos de prueba del experimento. Para reducir el sesgo de posición, el juez evalúa cada respuesta de forma absoluta, sin comparación con otra. Para limitar el sesgo de autopreferencia, el modelo juez (120B parámetros) difiere sustancialmente en escala de los modelos evaluados. Adicionalmente, las respuestas del juez se revisan manualmente para contrastar la fiabilidad del criterio automático, reconociendo que el juicio del LLM es una aproximación escalable al juicio humano, no un sustituto del mismo.

La evaluación humana sigue considerándose el estándar de referencia en la valoración de sistemas generativos, dado que los evaluadores humanos pueden apreciar matices semánticos, coherencia argumentativa y adecuación contextual que las métricas automáticas no capturan \cite{celikyilmaz2020evaluation} o que los LLMs no pueden replicar completamente. Su principal limitación es el coste y la dificultad de escalarla, lo que en la práctica obliga a combinarla con métodos automáticos como el \textit{LLM-as-a-judge}.